% Sections 8 and 9: Discussion and Concluding Remarks (new in A3).
\section{Discussion}\label{sec:discussion}

\subsection{Faster versus Greener, Then and Now (G1)}
Pinto et al.\ found the sequential version the most frugal in six of nine benchmarks and concluded that being faster had little correlation with being greener~\cite{pinto2014energy}. On both of our platforms the answer to RQ1.2 is the opposite: on CPU-bound work, parallel execution used 2.4 (Mac) to 4.0 (Intel) times less energy than sequential execution. On the M4, CPU power rose only from 3.9\,W with one busy thread to 10.1--11.2\,W with ten (AX4), while throughput grew near-linearly, so shared power is spread over more work. This is the race-to-idle regime that Kim et al.\ analyze~\cite{kim2015racing}, and on the Mac it makes energy track time ($\rho = 0.987$), as van Kempen et al.\ observed when power is held near constant~\cite{vankempen2025green}.

Intel shows the limit of that rule (RQ1.3, Figure~\ref{fig:evt}). The configurations that create OS threads continuously, PT and the per-operation pools at $p = 256$, ran at the 25\,W package limit with clocks about 1.8 times higher, finished up to 33\% faster, and used up to 38\% more energy. Speed bought with parallelism at a constant clock reduced energy; speed bought with frequency, and hence voltage, increased it, as in the compiler-flag case of Arag{\'o}n-Jurado et al.~\cite{aragonjurado2025faster}. We hypothesize (\textbf{H-a}) that frequent thread creation and wake-ups lead the hardware-managed P-state controller (HWP) to raise frequency, although the work itself does not need it. ``Faster is greener'' is thus a property of how a platform turns load into speed, not a general law.

\subsection{What Virtual Threads Actually Save (G2)}
At equal concurrency on blocking workloads, VT used 15 to 36 times less CPU energy than the platform-thread constructs on the Mac (RQ2.1). The control experiment shows that most of this gap is not the cost of waiting but the cost of creating threads: our A2 design, like the 2014 one, created FP and FJ inside every operation, whereas VT's carriers persist. Reused, the fixed pool matched VT on SLEEP (0.86 times) and beat it on Mac SOCKET (0.64 times). VT's remaining advantage is that a pool must be as wide as the number of concurrent waits, whereas VT needs only a few carriers. VT with few carriers was the lowest-energy configuration on both blocking workloads on Intel (4 carriers) and on Mac SLEEP (1 to 64 carriers). Conversely, 256 carriers is VT's worst setting, drawing 11 to 13\,W on Intel; we hypothesize (\textbf{H-b}) that carriers beyond the core count spend energy parking and unparking. Since the JDK's default carrier count equals the number of cores, the default configuration of VT is close to its efficient one. On MIXED, the most realistic kernel, energy is nearly equal and VT is the fastest on both platforms. As for RQ2.2, VT was significantly the lowest at only one level ($p = 256$ on CPU), against pools created per operation. Thread-per-task execution cannot reuse threads and stays expensive everywhere.

\subsection{Placement on Asymmetric Cores (G3)}
The cluster split (RQ3.2) shows that macOS places PT's threads mostly on P cores (E-cluster share 13 to 16\%) while VT's carriers stay on the E cores on SLEEP (100\%). Part of PT's 62- to 67-fold gap to VT on Mac SLEEP is therefore placement: confined to the E cores, it shrank to 3.2 to 3.4, close to Intel's 3.5 on two E cores. Pinning isolates core type: the same PT workload cost 1.5 times more on two P cores than on two E cores at the same speed, whereas VT cost the same on every class. On CPU work, E cores saved 5\% energy at 25\% more time, and LP-E cores lost on both counts as package static power accrued. The Mac's 8- to 15-fold saving under BACKGROUND (RQ3.1) is thus mostly a clock effect (1090 instead of 2752\,MHz) on a counter that excludes static power, while timer coalescing made sleeping work 7 to 11 times more expensive: E-core confinement suits CPU-bound batch work, not latency-bound work. We hypothesize (\textbf{H-c}) that the scheduler places threads by recent utilization: short-lived PT threads start on P cores, while mostly idle carriers settle on E cores.

\subsection{Measurement Scope Changes the Answer}
The Mac counter covers the CPU clusters only, so slow configurations look cheap: bounded pools with 8 to 64 workers used less CPU energy than VT on SOCKET while taking 3 to 20 times longer. Intel's package counter includes about 3.3\,W of idle and uncore power, so slow configurations pay for waiting, and idle-subtracted energy removes static power entirely. A ranking can invert between counters, so we report time with energy.

\subsection{Implications}
\emph{Java developers.} On CPU-bound work the construct hardly matters: use as many workers or carriers as cores. On blocking work, avoid creating OS threads per task or per request: either reuse a pool sized to the number of concurrent waits, or use virtual threads with the default carrier count, which matches a well-sized pool in energy without sizing and was the fastest on mixed work. \emph{JDK and OS developers.} Thread churn alone pushed the Intel package to its power limit, and keeping blocking work on efficiency cores, as macOS did for virtual threads, saves energy without code changes. \emph{Researchers.} Setup work inside the measured operation, such as creating a pool, can dominate a result, as our control showed.

\section{Concluding Remarks}\label{sec:conclusion}
On JDK~25 and two hybrid processors, parallel execution beat sequential execution on CPU-bound work, the opposite of the 2014 baseline; faster was greener on the M4 but not when Intel gained speed through frequency; and the blocking-workload advantage of virtual threads came mostly from not creating OS threads, which reused pools also achieve, while virtual threads need only a few carriers. Thread-per-task execution is costly on both platforms, partly because it runs on performance cores. \emph{Future work:} the full Intel campaign at 5 forks; experiments for H-a to H-c (capped clocks, carrier-count sweeps with JFR, scheduler traces); a wall-power meter for static power; and real services such as Spring PetClinic~\cite{stocker2026energyupgrades} with pools sized and reused as in production.
